\begin{table}[!htbp]
    \centering
    \caption{Pregled primerjanih modelov in njihovih vhodnih predstavitev.}
    \label{tab:pregled-primerjanih-modelov}
    \footnotesize
    \setlength{\aboverulesep}{0.25ex}
    \setlength{\belowrulesep}{0.4ex}
    \begin{tabularx}{\textwidth}{@{}>{\raggedright\arraybackslash}m{0.22\textwidth} >{\raggedright\arraybackslash}m{0.23\textwidth} >{\raggedright\arraybackslash}m{0.20\textwidth} >{\raggedright\arraybackslash}X@{}}
        \toprule
        \textbf{Družina modelov} & \textbf{Model} & \textbf{Vhod v model} & \textbf{Vloga v primerjavi} \\
        \midrule
        \multirow[c]{6}{=}{\raggedright izhodiščna klasifikatorja} & Naključni & \multirow[c]{6}{=}{\raggedright brez učenja iz značilk} & spodnja meja naključnega ugibanja \\
        \cmidrule(lr){2-2}\cmidrule(l){4-4}
        & Večinski & & spodnja meja napovedovanja najpogostejšega razreda \\
        \cmidrule{1-4}
        \multirow[c]{3.6}{=}{\raggedright klasični modeli strojnega učenja} & \ModelName{SVM} & \multirow[c]{3}{=}{\raggedright agregirane značilke okna} & \multirow[c]{3}{=}{\raggedright klasični model strojnega učenja} \\
        \cmidrule(lr){2-2}
        & \ModelName{LightGBM} & & \\
        \cmidrule(lr){2-2}
        & \ModelName{MLP} & & \\
        \cmidrule{1-4}
        \multirow[c]{6}{=}{\raggedright zamrznjeni prednaučeni predstavitveni modeli} & \ModelName{GazeMAE} & koordinati pogleda $(x, y)$ & prednaučena predstavitev za množico \emph{samo pogled} \\
        \cmidrule(l){2-4}
        & \ModelName{MOMENT} & časovne vrste izbranih signalov & prednaučena predstavitev za množice s signali zenic oziroma več signali \\
        \cmidrule{1-4}
        \multirow[c]{10.5}{=}{\raggedright grafovski modeli} & \ModelName{GCN} & homogeni graf okna & osnovni grafovski model \\
        \cmidrule(l){2-4}
        & \ModelName{HeteroGCN-mean} & \multirow[c]{3.5}{=}{\raggedright heterogeni graf okna} & ločena obravnava relacij s preprosto fuzijo \\
        \cmidrule(lr){2-2}\cmidrule(l){4-4}
        & \ModelName{HeteroGCN-MLP} & & ločena obravnava relacij z učljivo fuzijo \\
        \cmidrule(l){2-4}
        & \ModelName{HeteroGCN-MLP-w} & heterogeni graf okna in značilke povezav & najkompleksnejša stopnja arhitekturne lestvice \\
        \bottomrule
    \end{tabularx}
\end{table}
